\section{Quantum-Secure Settlement}
\label{sec:quantum}

\subsection{Three independent hardness assumptions}

\Quasar{} 3.0 (\Lux{} \textsc{lp-020}, \cite{lp-020-quasar-3,lp-105-quasar})
provides finality based on \emph{three} independent cryptographic hardness
assumptions, composed in a single block-finalization certificate. An
adversary must break \emph{all three} simultaneously to forge a
\Lux{}/\Zoo{} block.

\begin{longtable}{p{3.8cm}p{4.5cm}p{6.4cm}}
\toprule
\textbf{Layer} & \textbf{Hardness assumption} & \textbf{Role in finality} \\
\midrule
\endhead
$\BLS{}_{12\text{-}381}$ & Discrete log over BLS pairings (classical).
   See \Lux{} \textsc{lp-075}~\cite{lp-075-bls}.
   & Classical aggregate signature for the validator quorum.
   Carries the bulk of liveness throughput. \\
Ringtail Ring-LWE & Ring Learning-With-Errors (post-quantum lattice).
   See \Lux{} \textsc{lp-073}~\cite{lp-073-ringtail,ring-lwe}.
   & Post-quantum threshold signature; co-signed with \BLS{} on every
   block. Falls back to majority Ringtail-only certificate if \BLS{}
   is suspect. \\
$\MLDSA{}\text{-}65$ via Z-Chain Groth16 &
   Module-Lattice signature (\textsc{nist} \textsc{fips}~204).
   See \Lux{} \textsc{lp-070}~\cite{lp-070-mldsa,nist-mldsa}, with
   succinct proof of validator identity via Groth16~\cite{groth16}.
   & Post-quantum \emph{identity proof} layer; binds each
   block-signing key to a \textsc{pq}-safe credential committed on
   Z-Chain. \\
\bottomrule
\end{longtable}

\subsection{Forgery requires breaking all three}

A valid \Quasar{} 3.0 block-finalization certificate $\sigma$ contains:

\[
\sigma \;=\; \big( \sigma_{\BLS{}},\; \sigma_{\text{Ringtail}},\; \pi_{\MLDSA{}} \big)
\]

\noindent where $\sigma_{\BLS{}}$ is the aggregate $\BLS{}_{12\text{-}381}$
signature from the validator quorum, $\sigma_{\text{Ringtail}}$ is the
threshold Ring-LWE signature, and $\pi_{\MLDSA{}}$ is a Groth16 proof
that every signing key was bound to a valid $\MLDSA{}\text{-}65$
credential at block height $h$.

\begin{itemize}[nosep,leftmargin=1.2em]
\item Forging $\sigma_{\BLS{}}$ alone requires breaking discrete log on
  BLS12-381. Classical adversary; quantum adversary breaks this in
  polynomial time post-Q-day.
\item Forging $\sigma_{\text{Ringtail}}$ requires solving Ring-LWE for
  the chosen lattice parameters. Believed secure against both
  classical and quantum adversaries.
\item Forging $\pi_{\MLDSA{}}$ requires either breaking
  $\MLDSA{}\text{-}65$ (a different module-lattice problem from
  Ring-LWE) or breaking the Groth16 soundness assumption (knowledge
  of pairing-based extractor).
\end{itemize}

The hardness assumptions are deliberately heterogeneous: discrete log
over pairing-friendly curves (classical), Ring-LWE (lattice family A),
and module lattice signatures with succinct proof (lattice family B
with pairing-based zero-knowledge). A break in any single family does
not propagate to the others.

\subsection{Q-day resilience}

Define \emph{Q-day} as the day a cryptographically relevant quantum
computer breaks $\BLS{}_{12\text{-}381}$ in polynomial wall-clock time.

\begin{itemize}[nosep,leftmargin=1.2em]
\item \emph{Pre-Q-day}: all three signatures required. \BLS{} carries
  liveness; \textsc{pq} layers are belt-and-braces.
\item \emph{Post-Q-day}: \BLS{} compromise alone does not break
  finality. The Ringtail + $\MLDSA{}$ pair is sufficient for
  Byzantine-fault-tolerant finalization. The chain enters
  ``\textsc{pq}-only'' mode by social-consensus governance vote
  (Holographic \DAO{}, \S\ref{sec:dao}) until \BLS{} is replaced.
\item \emph{No silent break}: an attacker who silently obtains a
  \BLS{} forgery cannot retroactively rewrite history because the
  Ringtail + $\MLDSA{}$ proofs remain on-chain in every prior block.
\end{itemize}

\subsection{Settlement guarantees for Zoo / Liquidity workloads}

Every securities trade settled via the \Lux{} chain stack carries a
\Quasar{} 3.0 certificate. Three concrete guarantees follow for a
retail user:

\begin{enumerate}[nosep,leftmargin=1.4em]
\item \textbf{Trade integrity is post-quantum-safe.} The trade match
  is signed under all three layers; a future quantum adversary cannot
  rewrite a settled trade.
\item \textbf{Custody is post-quantum-safe.} M-Chain (\Lux{}
  \textsc{lp-019}~\cite{lp-019-mchain}) holds threshold custody keys
  using Ringtail \textsc{pq} threshold signatures.
\item \textbf{Audit trail is post-quantum-safe.} The Merkle root of
  every settled block is anchored under \Quasar{} 3.0; the auditor
  needs only the \textsc{pq} verification key set to verify
  historical state.
\end{enumerate}

\subsection{Cross-references}

\begin{itemize}[nosep,leftmargin=1.2em]
\item \Lux{} \textsc{lp-020} --- \Quasar{} Consensus 3.0~\cite{lp-020-quasar-3}.
\item \Lux{} \textsc{lp-070} --- $\MLDSA{}\text{-}65$ profile~\cite{lp-070-mldsa}.
\item \Lux{} \textsc{lp-073} --- Ringtail Ring-LWE threshold~\cite{lp-073-ringtail}.
\item \Lux{} \textsc{lp-075} --- $\BLS{}_{12\text{-}381}$ aggregate~\cite{lp-075-bls}.
\item \Lux{} \textsc{lp-105} --- \Quasar{} 3.0 paper~\cite{lp-105-quasar}.
\item \Lux{} \textsc{lp-019} --- M-Chain MPC~\cite{lp-019-mchain}.
\item \Lux{} \textsc{lp-013} --- F-Chain FHE~\cite{lp-013-fchain}.
\item \Lux{} \textsc{lp-134} --- A/B/M/F chain topology~\cite{lp-134-topology}.
\end{itemize}
